\documentclass{article}
\usepackage{lmodern}
\usepackage{amsmath}
\usepackage{amssymb}
\usepackage{xfrac}
\usepackage{cancel}
\usepackage{accents}
\providecommand{\utilde}[1]{\underaccent{\tilde}{#1}}
\usepackage[letterpaper,margin=1in]{geometry}

\title{\textbf{TeX Live \texttt{mlist.c} Parity Verification Suite}}
\author{RushTeX Differential Testing Framework}
\date{September 2026}

\begin{document}

\maketitle

This suite exercises the 56 mathematical typesetting routines and TeX Appendix~G rules from TeX Live's \texttt{mlist.c}.

\section{Atom Classes \& TeX Appendix G Rule 5 Normalization}
The 8 math atom classes (Ord, Op, Bin, Rel, Open, Close, Punct, Inner) and inter-atom spacing:
\[
x + y \cdot z = \sin(\theta) \pm \sqrt{2}, \quad a, b \in \mathbb{R}
\]
Rule 5 converts binary operators to ordinary atoms when preceded or followed by non-compatible atoms:
\begin{itemize}
    \item Leading binary sign: $+x$ and $-y$ (converted to Ord).
    \item Trailing and fence-adjacent operators: $(+x)$, $[*a]$, and $(x+)$.
    \item Consecutive relations and operators: $x + = y$ and $a = = b$.
\end{itemize}

\section{Fractions and Generalized Quotients (Appendix G Rule 15)}
Display and inline fraction layouts with rule thickness and vertical clearance:
\[
\frac{a + b}{c + d} = 1 + \frac{x^2}{1 + \frac{x^3}{1 + x^4}}, \quad \binom{n}{k} = \frac{n!}{k!(n-k)!}
\]
Inline fractions and slashed forms: $\frac{1}{2}$, $\frac{x+y}{x-y}$, and $\sfrac{3}{4}$.

\section{Radicals and Root Indices (Appendix G Rule 11)}
Square roots of varying heights and multi-level root nesting:
\[
\sqrt{x} < \sqrt{x + y} < \sqrt{x^2 + y^2} < \sqrt{\frac{a^2 + b^2}{c^2 + d^2}} < \sqrt{1 + \sqrt{2 + \sqrt{3 + \sqrt{4}}}}
\]
Roots with arbitrary degree index: $\sqrt[3]{x}$, $\sqrt[n]{1 + x^n}$, $\sqrt[\alpha + \beta]{e^x}$.

\section{Large Operators and Limit Placements (Appendix G Rule 13)}
Summations, products, integrals, and unions with display vs text limits:
\[
\sum_{i=1}^n i^3 = \left( \sum_{i=1}^n i \right)^2, \quad \prod_{k=1}^\infty \left( 1 - \frac{1}{p_k^2} \right) = \frac{6}{\pi^2}
\]
Integrals with differential spacing:
\[
\int_0^\infty e^{-x^2} \, \mathrm{d}x = \frac{\sqrt{\pi}}{2}, \quad \oint_C \vec{F} \cdot \mathrm{d}\vec{r} = \iint_S (\nabla \times \vec{F}) \cdot \mathrm{d}\vec{S}
\]
Explicit limit overrides: $\sum\nolimits_{i=1}^n x_i$ and $\int\limits_a^b f(t)\,\mathrm{d}t$.

\section{Scripts and TeX Rule 12 Accent Attachment (Appendix G Rule 18)}
Subscripts, superscripts, multi-level towers, and complex drops:
\[
x^2, \quad x_i, \quad x_i^2, \quad A_{i,j,k}^{2\prime}, \quad 2^{2^{2^x}}, \quad \Gamma_0^\dagger, \quad f_1(x)
\]
Rule 12 Script Attachment on Accents (attaching scripts relative to the base nucleus):
\[
\hat{x}^2, \quad \vec{v}_0^2, \quad \tilde{a}_i^\prime, \quad \bar{X}_{n+1}^k, \quad \dot{r}^2 + r^2 \dot{\theta}^2
\]

\section{Mathematical Accents: Top, Bottom, and Overlay}
Single and wide accents placed with skewchar kerning:
\[
\hat{a}, \quad \tilde{b}, \quad \bar{c}, \quad \vec{v}, \quad \dot{x}, \quad \ddot{y}, \quad \breve{z}, \quad \check{w}, \quad \acute{u}, \quad \grave{e}
\]
Wide accents spanning composite nucleus expressions:
\[
\widehat{x + y}, \quad \widetilde{A + B + C}, \quad \overline{z_1 + z_2}, \quad \underline{u \cdot v}
\]
Bottom and overlay accents:
\[
\utilde{x}, \quad x \not= y, \quad \cancel{k}
\]

\section{Delimiters and Extensible Structures (Appendix G Rule 19)}
Fixed size delimiter tiers:
\[
(x) \quad \bigl( x \bigr) \quad \Bigl( x \Bigr) \quad \biggl( x \biggr) \quad \Biggl( x \Biggr)
\]
Dynamically stretched fences enclosing tall inner expressions:
\[
\left( \frac{x^2 + 1}{x^2 - 1} \right), \quad \left[ 1 + \left( \frac{a}{b} \right)^3 \right], \quad \left\{ x \in \mathbb{R} \;\middle|\; \sqrt{x} > 1 \right\}
\]
Over- and under-delimiters with centered labels:
\[
\overbrace{x_1 + x_2 + \dots + x_n}^{n \text{ elements}} = \underbrace{y_1 + y_2 + \dots + y_m}_{m \text{ elements}}
\]
Extensible horizontal arrows:
\[
A \xrightarrow{f} B \xleftarrow{g \circ h} C
\]

\section{Spacing Tables and Conditional Math Glue}
Explicit spacing commands:
\[
x \, y \quad x \: y \quad x \; y \quad x \quad y \quad x \qquad y \quad x \! y
\]
Conditional math glue (\texttt{\textbackslash nonscript}):
\[
a \nonscript\, b \quad \text{vs.} \quad x_{a \nonscript\, b}
\]

\section{Inline Math Paragraph Line Breaking (mlist.c:4757--4789)}
In this paragraph, we verify that inline mathematical formulas legally break across lines at binary operators ($\binoppenalty = 700$) and relational operators ($\relpenalty = 500$):
Let us examine the following continuous algebraic identity placed inside running body text to inspect break behavior: $T(x) = a_0 + a_1 x + a_2 x^2 + a_3 x^3 + a_4 x^4 + a_5 x^5 = \sum_{k=0}^5 a_k x^k = \Phi(x) + \Psi(x) = \Omega(x)$, where each operator provides a candidate breakpoint with proper glue discardability at the line wrap.

\end{document}
